\documentclass{article}
\usepackage{hyperref}

\begin{document}
    \begin{center}
        \Large \textbf{Challange 1: Quick Start} \\ [6pt]
        \normalsize
        Autor: Maximilian Jakob Maag B.Sc. \\
        Version: 1.0
    \end{center}
    \section{Intro}
    Dieses Handout gibt Hinweise zum Setup und soll Dokumente und Informationen für Sie zusammentragen.
    \section{Empfohlenes Setup}
    Empfohlen wird die Verwendung von Linux auf einer Debian oder Arch basierenden Distro mit der ZSH oder BASH Shell. \\
    Eine Teilnahme unter Windows ist ebenfalls möglich wurde aber nicht getestet. \\ \\
    Es wird dringend empfohlen sich auf eine Shell festzulegen. Sensible Informationen, z.B. \texttt{API\_TOKEN} sollten in Umgebungsvariablen gespeichert werden.
    Passwörter und private keys dürfen sich nicht in Git Repositories befinden. Ein entsprechendes Git Ignore file sollte angelegt werden.
    \section{Dokumentation und Wikis}
    Nachstehend eine Liste mit Wikis und Dokumentationen. \\ \\
    \begin{itemize}
        \item \href{https://git-scm.com/book/en/v2}{Pro Git} Ein Buch über Git und dessen Verwendung.
        \item \href{https://developer.hashicorp.com/terraform/docs}{Terraform Doku} Dokumentation über Terraform vom Hersteller.
        \item \href{https://registry.terraform.io/providers/linode/linode/latest/docs}{Linode Provider} Schnittstellen Doku für Linode. Andere Provider nutzen eine eigene Doku.
    \end{itemize}
\end{document}